\documentclass[a4paper]{article}

\usepackage{pgfplots}
\pgfplotsset{compat=1.18}
\usepackage[T1]{fontenc}
\usepackage[utf8]{inputenc}
\usepackage[italian]{babel}
\usepackage{amsfonts}
\usepackage{amsmath}
\usepackage{amssymb}
\usepackage{amsthm}
\usepackage{newtxtext}
\usepackage{newtxmath}
\usepackage{tikz}

\usepackage[margin=1.5cm]{geometry}
\usepackage{hyperref}
\usetikzlibrary{arrows.meta, patterns, decorations.markings, decorations.pathreplacing}

\tikzset{
    ->-/.style={decoration={markings, mark=at position #1 with \arrow{Stealth[length=2.5mm]}}, postaction={decorate}}
}

\setlength{\parindent}{0pt}

\theoremstyle{plain}
\newtheorem{teorema}{Teorema}

\theoremstyle{definition}
\newtheorem{definizione}{Definizione}

\theoremstyle{remark}
\newtheorem{osservazione}{Osservazione}

\newcommand{\N}{\mathbb{N}}
\newcommand{\Z}{\mathbb{Z}}
\newcommand{\Q}{\mathbb{Q}}
\newcommand{\R}{\mathbb{R}}
\newcommand{\set}[1]{\{#1\}}

\title{Appunti di Matematica}
\author{Marco Galanti}
\date{\today}

\begin{document}

\maketitle
\tableofcontents
\newpage
\section{Equazioni e disequazioni di grado superiore al secondo}
\subsection{Equazioni monomie, binomie e trinomie}
\begin{definizione}
    Un'equazione nell'incognita $x$ si dice monomia se è riconducibile alla forma
    \begin{equation}
        ax^n=0
    \end{equation}
    dove $a\in\R\setminus\set{0}$ ed $n\in\Z^+$.
\end{definizione}
Esempio: $3x^2=0\implies x=0$
\begin{definizione}
    Un'equazione nell'incognita $x$ si dice binomia se è riconducibile alla forma:
    \begin{equation}
        ax^n+b=0
    \end{equation}
    dove $a,b\in\R\setminus\set{0}$ ed $n\in\Z^+$.
\end{definizione}
L'equazione si riconduce a $x^n=-\dfrac{b}{a}$.

Esempio 1: $(x+10)^3=-8$

$\begin{array}{l}
    (x+10)^3=-8\\
    x+10=\sqrt[3]{-8}\\
    x+10=-2\\
    x=-12
\end{array}$

Esempio 2: $4(x-1)^4=25$

$\begin{array}{l}
    4(x-1)^4=25\\
    (x-1)^4=\dfrac{25}{4}\\
    x-1=\pm\sqrt[4]{\dfrac{25}{4}}\\
    x=1\pm\sqrt{\dfrac{5}{2}}=1\pm\dfrac{\sqrt{10}}{2}
\end{array}$
\begin{definizione}
    Un'equazione nell'incognita $x$ si dice trinomia se è riconducibile alla forma:
    \begin{equation}
        ax^{2n}+bx^n+c=0
    \end{equation}
    dove $a,b,c\in\R\setminus\set{0}$ ed $n\in\Z^+$.
\end{definizione}
Se $n=1$ l'equazione trinomia è riconducibile a un'equazione di $\mathrm{II}^\circ$ grado.

Se $n=2$ l'equazione trinomia è nella forma $ax^4+bx^2+c=0$ ed è detta biquadratica.

Per $n\geq2$ si pone $x^n=t$, riconducendo l'equazione trinomia in $at^2+bt+c=0$.

L'equazione si risolve in $t$ e poi si pone $x^n=t_1$ e $x^n=t_2$ per risolvere l'equazione in $x$.
\newpage
Esempio 1: $x^4+2x^2-3=0$

$\begin{array}{l}
   x^4+2x^2-3=0\quad x^2=t\\
   t^2+2t-3=0\\
   t=-3\lor t=1\\
   x^2=-3\lor x^2=1\\
   \nexists x\in\R\lor x=\pm1\\
   x=\pm1
\end{array}$

Esempio 2: $x^4+2x^2+3=0$

$\begin{array}{l}
    x^4+2x^2+3=0\quad x^2=t\\
    t^2+2t+3=0\\
    \Delta<0\implies\nexists t\in\R\\
    \nexists x\in\R
\end{array}$

Esempio 3: $x^6+3x^3+2=0$

$\begin{array}{l}
    x^6+3x^3+2=0\quad x^3=t\\
    t^2+3t+2=0\\
    t=-2\lor t=-1\\
    x^3=-2\lor x^3=-1\\
    x=\sqrt[3]{-2}=-\sqrt[3]{2}\lor x=\sqrt[3]{-1}=-1\\
    x=-\sqrt[3]{2}\lor x=-1
\end{array}$

Esempio 4: $x^8-5x^4+4=0$

$\begin{array}{l}
   x^8-5x^4+4=0\quad x^4=t\\
   t^2-5t+4=0\\
   t=1\lor t=4\\
   x^4=1\lor x^4=4\\
   x=\pm\sqrt[4]{1}=\pm1\lor x=\pm\sqrt[4]{4}=\pm\sqrt{2}\\
   x=\pm1\lor x=\pm\sqrt{2}
\end{array}$
\newpage
\subsection{Equazioni risolvibili mediante scomposizione in fattori}
Un'equazione nella forma $P(x)=0$, dove $P(x)$ è un polinomio nella variabile $x$, può essere risolta ponendo ogni fattore uguale a $0$ sfruttando la legge di annullamento del prodotto.

Esempio 1: $x^3+8=0$

$\begin{array}{l}
    x^3+8=0\\
    (x+2)(x^2-2x+4)=0\\
    x+2=0\lor x^2-2x+4=0\\
    x=-2\lor\nexists x\in\R\quad(\Delta<0)\\
    x=-2
\end{array}$

Esempio 2: $x^4-3x^2+2=0$

$\begin{array}{l}
    x^4-3x^2+2=0\\
    (x^2-1)(x^2-2)=0\\
    x^2-1=0\lor x^2-2=0\\
    x=\pm1\lor x=\pm\sqrt{2}
\end{array}$

\newpage
\subsection{Disequazioni intere di grado superiore al secondo}
Si chiamano monomie le disequazioni nelle seguenti forme, con $a\ne0$ ed $n\in\N\setminus\set{0}$:
\begin{equation}
    ax^n<0\quad ax^n>0\quad ax^n\leq0\quad ax^n\geq0
\end{equation}

Esempio 1: $3x^4>0$

Il primo membro è sempre positivo, eccetto che per il valore $x=0$ per cui si annulla.

Di conseguenza la disequazione è verificata $\forall x\in\R\setminus\set{0}$.

Esempio 2: $8x^5\leq0$

Poiché una potenza di esponente dispari ha sempre lo stesso segno della base, la disequazione è verificata per $x\leq0$.

Si chiamano binomie le disequazioni nelle seguenti forme, con $a\neq0$, $b\neq0$ ed $n\in\N\setminus\set{0}$:
\begin{equation}
    ax^n+b<0\quad ax^n+b>0\quad ax^n+b\leq0\quad ax^n+b\geq0
\end{equation}

Formule risolutive di $ax^n+b>0$ con $a>0$ e $b\ne0$:
\begin{enumerate}
    \item $n$ dispari: $x>\sqrt[n]{-\dfrac{b}{a}}$
    \item $n$ pari:
    \begin{enumerate}
        \item $b>0$: $\forall x\in\R$
        \item $b<0$: $x<-\sqrt[n]{-\dfrac{b}{a}}\lor x>\sqrt[n]{-\dfrac{b}{a}}$
    \end{enumerate}
\end{enumerate}
Formule risolutive di $ax^n+b<0$ con $a>0$ e $b\ne0$:
\begin{enumerate}
    \item $n$ dispari: $x<\sqrt[n]{-\dfrac{b}{a}}$
    \item $n$ pari:
    \begin{enumerate}
        \item $b>0$: $\nexists x\in\R$
        \item $b<0$: $-\sqrt[n]{-\dfrac{b}{a}}<x<\sqrt[n]{-\dfrac{b}{a}}$
    \end{enumerate}
\end{enumerate}
Esempio 1: $3x^5+1>0\implies 3x^5>-1\implies x^5>-\dfrac{1}{3}\implies x>\sqrt[5]{-\dfrac{1}{3}}\implies x>-\dfrac{1}{\sqrt[5]{3}}$

Esempio 2: $-8x^3+1>0\implies -8x^3>-1\implies x^3<\dfrac{1}{8}\implies x<\sqrt[3]{\dfrac{1}{8}}\implies x<\dfrac{1}{2}$

Si chiamano trinomie le disequazioni nelle seguenti forme:
\begin{equation}
    ax^{2n}+bx^n+c<0\quad ax^{2n}+bx^n+c>0\quad ax^{2n}+bx^n+c\leq0\quad ax^{2n}+bx^n+c\geq0
\end{equation}

Per risolvere queste disequazioni si pone $x^n=t$ e ci si riconduce a una disequazione di $\mathrm{II}^\circ$ grado.

Esempio 1: $x^4-3x^2+4\leq0$

$\begin{array}{l}
   x^2=t\quad x^4=t^2\\
   t^2-3t+4\leq0\\
   \Delta<0\land a>0\implies t^2-3t+4>0\ \ \forall t\in\R\\
   \text{diseq. }\nexists x\in\R
\end{array}$

Esempio 2: $x^6-3x^3+2\geq0$

$\begin{array}{l}
   x^3=t\quad x^6=t^2 \\
   t^2-3t+2\geq0\\
   t\leq1\lor t\geq2\\
   x^3\leq1\lor x^3\geq2\\
   x\leq1\lor x\geq\sqrt[3]{2}
\end{array}$

Per la risoluzione di una disequazione nelle forme, dove $P(x)$ è un polinomio nella variabile $x$:
\begin{equation}
    P(x)>0\quad P(x)<0\quad P(x)\geq0\quad P(x)\leq0
\end{equation}
si può ricorrere alla scomposizione in fattori e allo sfruttamento della legge di annullamento del prodotto.

Esempio: $x^3-x^2\leq 4x-4$

$\begin{array}{l}
    x^3-x^2-4x+4\leq0\\
    x^2(x-1)-4(x-1)\leq0\\
    (x-1)(x^2-4)\leq0\\
    F_1>0\quad x-1>0\implies x>1\\
    F_2>0\quad x^2-4>0\implies x<-2\lor x>2
\end{array}$

\begin{tikzpicture}[x=0.5cm,y=0.5cm,font=\footnotesize]

  \def\xa{1.5}
  \def\xb{5.5}
  \def\xc{8.5}

  \draw[-{Stealth[length=2mm]}, thick] (-0.5,0) -- (11,0) node[below left] {$x$};

  \node[above] at (\xa,0.15) {$-2$};
  \node[above] at (\xb,0.15) {$1$};
  \node[above] at (\xc,0.15) {$2$};

  \foreach \x in {\xa,\xb,\xc}
    \draw[thick] (\x,0.2) -- (\x,-0.2);

  \draw (-0.5,-1) -- (10,-1);
  \draw (-0.5,-2) -- (10,-2);
  \draw (-0.5,-3) -- (10,-3);

  \draw[very thick] (\xa,0) -- (\xa,-1);
  \draw[very thick] (\xc,0) -- (\xc,-1);
  \draw[very thick] (\xb,-1) -- (\xb,-2);

  \foreach \x in {\xa,\xb,\xc}
    \draw[very thick] (\x,-3) -- (\x,-3.5);

  \node at (0.5,-0.5) {$-$};
  \node at (3.5,-0.5) {$-$};
  \node at (\xb,-0.5) {$0$};
  \node at (7,-0.5)   {$+$};
  \node at (9.5,-0.5) {$+$};

  \node at (0.5,-1.5) {$+$};
  \node at (\xa,-1.5) {$0$};
  \node at (3.5,-1.5) {$-$};
  \node at (7,-1.5)   {$-$};
  \node at (\xc,-1.5) {$0$};
  \node at (9.5,-1.5) {$+$};

  \node at (0.5,-2.5) {$-$};
  \node at (\xa,-2.5) {$0$};
  \node at (3.5,-2.5) {$+$};
  \node at (\xb,-2.5) {$0$};
  \node at (7,-2.5)   {$-$};
  \node at (\xc,-2.5) {$0$};
  \node at (9.5,-2.5) {$+$};

\end{tikzpicture}

$x\leq-2\lor1\leq x\leq2$

Si chiamano fratte le disequazioni nelle seguenti forme, dove $A(x)$ e $B(x)$ sono polinomi nella variabile $x$:
\begin{equation}
    \dfrac{A(x)}{B(x)}<0\quad \dfrac{A(x)}{B(x)}>0\quad \dfrac{A(x)}{B(x)}\leq0\quad \dfrac{A(x)}{B(x)}\geq0
\end{equation}

Esempio: $\dfrac{x-5}{4-x^2}\leq0$

$\begin{array}{l}
    N>0\quad x-5>0\implies x>5\\
    D>0\quad 4-x^2>0\implies -2<x<2
\end{array}$

\begin{tikzpicture}[x=0.5cm,y=0.5cm,font=\footnotesize,
  v/.style={fill=white,inner sep=1pt}]

  \def\xa{3}
  \def\xb{7}
  \def\xc{11}

  \draw[-{Stealth[length=2mm]}, thick] (-0.5,0) -- (14.5,0) node[below left] {$x$};

  \node[above] at (\xa,0.15) {$-2$};
  \node[above] at (\xb,0.15) {$2$};
  \node[above] at (\xc,0.15) {$5$};

  \draw (-0.5,-1) -- (13.5,-1);
  \draw (-0.5,-2) -- (13.5,-2);
  \draw (-0.5,-3) -- (13.5,-3);

  \foreach \x in {\xa,\xb,\xc}
    \draw[very thick] (\x,0.2) -- (\x,-3.5);

  \node at (1,-0.5)   {$-$};
  \node at (5,-0.5)   {$-$};
  \node at (9,-0.5)   {$-$};
  \node[v] at (\xc,-0.5) {$0$};
  \node at (13,-0.5)  {$+$};

  \node at (1,-1.5)   {$-$};
  \node[v] at (\xa,-1.5) {$0$};
  \node at (5,-1.5)   {$+$};
  \node[v] at (\xb,-1.5) {$0$};
  \node at (9,-1.5)   {$-$};
  \node at (13,-1.5)  {$-$};

  \node at (1,-2.5)   {$+$};
  \node[v] at (\xa,-2.5) {$\nexists$};
  \node at (5,-2.5)   {$-$};
  \node[v] at (\xb,-2.5) {$\nexists$};
  \node at (9,-2.5)   {$+$};
  \node[v] at (\xc,-2.5) {$0$};
  \node at (13,-2.5)  {$-$};

\end{tikzpicture}

$-2<x<2\lor x\geq5$
\begin{definizione}
    Un insieme di $2$ o più disequazioni, tutte nella stessa incognita, che si vuole siano soddisfatte contemporaneamente, si chiama sistema di disequazioni.
\end{definizione}
Esempio: $\begin{cases}
    \dfrac{5-x}{x+2}\geq0\\
    \dfrac{1}{2}x^2+x>4
\end{cases}$

$\begin{array}{l}
    N>0\quad 5-x>0\implies -x>-5\implies x<5\\
    D>0\quad x+2>0\implies x>-2
\end{array}$

\begin{tikzpicture}[x=0.5cm,y=0.5cm,font=\footnotesize,
  v/.style={fill=white,inner sep=1pt}]

  \def\xa{3}
  \def\xb{9}

  \draw[-{Stealth[length=2mm]}, thick] (-0.5,0) -- (13,0) node[below left] {$x$};

  \node[above] at (\xa,0.15) {$-2$};
  \node[above] at (\xb,0.15) {$5$};

  \draw (-0.5,-1) -- (12,-1);
  \draw (-0.5,-2) -- (12,-2);
  \draw (-0.5,-3) -- (12,-3);

  \foreach \x in {\xa,\xb}
    \draw[very thick] (\x,0.2) -- (\x,-3.5);

  \node at (1,-0.5)   {$+$};
  \node at (6,-0.5)   {$+$};
  \node[v] at (\xb,-0.5) {$0$};
  \node at (11,-0.5)  {$-$};

  \node at (1,-1.5)   {$-$};
  \node[v] at (\xa,-1.5) {$0$};
  \node at (6,-1.5)   {$+$};
  \node at (11,-1.5)  {$+$};

  \node at (1,-2.5)   {$-$};
  \node[v] at (\xa,-2.5) {$\nexists$};
  \node at (6,-2.5)   {$+$};
  \node[v] at (\xb,-2.5) {$0$};
  \node at (11,-2.5)  {$-$};

\end{tikzpicture}

$\begin{array}{l}
    -2<x\leq5\\
    \dfrac{1}{2}x^2+x>4\implies x^2+2x>8\implies x^2+2x-8>0\\
    x<-4\lor x>2
\end{array}$

\begin{tikzpicture}[x=0.5cm,y=0.5cm,font=\footnotesize,
  open/.style={draw,fill=white,circle,inner sep=0pt,minimum size=4pt},
  closed/.style={draw,fill=black,circle,inner sep=0pt,minimum size=4pt}]

  \def\xa{3}
  \def\xb{7}
  \def\xc{11}
  \def\xd{15}

  \fill[gray!20] (\xc,0) rectangle (\xd,-7);

  \draw[-{Stealth[length=2mm]}, thick] (-0.5,0) -- (19.5,0) node[below left] {$x$};

  \node[above] at (\xa,0.15) {$-4$};
  \node[above] at (\xb,0.15) {$-2$};
  \node[above] at (\xc,0.15) {$2$};
  \node[above] at (\xd,0.15) {$5$};

  \foreach \x in {\xa,\xb,\xc,\xd}
    \draw[very thick] (\x,0.2) -- (\x,-7);

  \draw[thick,dashed] (-0.5,-2) -- (\xb,-2);
  \draw[thick]        (\xb,-2) -- (\xd,-2);
  \draw[thick,dashed] (\xd,-2) -- (19,-2);
  \node[open]   at (\xb,-2) {};
  \node[closed] at (\xd,-2) {};

  \draw[thick]        (-0.5,-4) -- (\xa,-4);
  \draw[thick,dashed] (\xa,-4) -- (\xc,-4);
  \draw[thick]        (\xc,-4) -- (19,-4);
  \node[open] at (\xa,-4) {};
  \node[open] at (\xc,-4) {};

  \draw[thick,dashed] (-0.5,-6) -- (\xc,-6);
  \draw[thick]        (\xc,-6) -- (\xd,-6);
  \draw[thick,dashed] (\xd,-6) -- (19,-6);
  \node[open]   at (\xc,-6) {};
  \node[closed] at (\xd,-6) {};

\end{tikzpicture}

$2<x\leq5$
\newpage
\section{Equazioni e disequazioni irrazionali}
\begin{definizione}
    Un'equazione o una disequazione si chiama irrazionale quando contiene almeno un radicale nel cui radicando compare l'incognita.
\end{definizione}
Se l'indice del radicale è dispari, sia in un'equazione sia in una disequazione, elevando all'indice entrambi i membri si ottiene un'equazione, o una disequazione, equivalente.
Se invece l'indice è pari allora si procede nel seguente modo, chiamato $n$ l'indice, con $n\in\N$ e $n\geq2$:
\begin{equation}
    \sqrt[n]{A(x)}=B(x)\iff\begin{cases}
        B(x)\geq0\\
        A(x)=B^n(x)
    \end{cases}
\end{equation}
In questo sistema è esclusa la condizione di esistenza del radicale, $A(x)\geq0$, poiché è già soddisfatta nella seconda equazione, dovendo $A(x)$ essere uguale a una potenza di esponente pari. $B(x)\ge0$ è la concordanza del segno, dovendo un radicale con indice pari essere uguale a una quantità non negativa.

Esempio: $\sqrt{x^2+4}=2x-2$

$\begin{array}{l}
    \begin{cases}
        2x-2\geq0\\
        x^2+4=(2x-2)^2
    \end{cases}\\
    \begin{cases}
        x\geq1\\
        3x^2-8x=0
    \end{cases}\\
    \begin{cases}
        x\geq1\\
        x=0\lor x=\dfrac{8}{3}
    \end{cases} \\
    x=\dfrac{8}{3}
\end{array}$

Se il secondo membro è un numero $k\in\R$ allora esistono delle tecniche per sveltire i calcoli:
\begin{enumerate}
    \item $k<0$: $\nexists x\in\R$
    \item $k\geq0$: $A(x)=k^n$
\end{enumerate}
Se in un'equazione compare più di un radicale si procede nel seguente modo:
\begin{enumerate}
    \item Si determina il dominio dell'equazione ponendo ogni radicando non negativo e ogni denominatore non nullo;
    \item Dopo aver isolato i radicali in modo che entrambi i membri siano non negativi, si elevano entrambi i membri all'indice comune dei radicali che vi compaiono.
\end{enumerate}
Esempio: $\sqrt{x}-\sqrt{x+5}=-\sqrt{x-3}$

Condizioni di esistenza: $\begin{cases}
    x\geq0\\
    x+5\geq0\\
    x-3\geq0
\end{cases}\implies x\geq3$

$\begin{array}{l}
    \sqrt{x}+\sqrt{x-3}=\sqrt{x+5}\\
    \left(\sqrt{x}+\sqrt{x-3}\right)^2=\left(\sqrt{x+5}\right)^2\\
    x+x-3+2\sqrt{x^2-3x}=x+5\\
    2\sqrt{x^2-3x}=8-x\\
    \begin{cases}
        8-x\geq0\\
        4(x^2-3x)=(8-x)^2
    \end{cases}\\
    x=-\dfrac{16}{3}\lor x=4\\
    \begin{cases}
        3\leq x\leq8\\
        x=-\dfrac{16}{3}\lor x=4
    \end{cases} \implies x=4
\end{array}$
\newpage
Per risolvere disequazioni irrazionali nella forma $\sqrt[n]{A(x)}<B(x)$ con $n$ pari si riconducono al seguente metodo:
\begin{equation}
    \sqrt[n]{A(x)}<B(x)\iff\begin{cases}
        A(x)\geq0\\
        B(x)>0\\
        A(x)<B^n(x)
    \end{cases}
\end{equation}
In caso venisse presa in considerazione anche l'uguaglianza si ricorre al seguente metodo:
\begin{equation}
    \sqrt[n]{A(x)}\leq B(x)\iff\begin{cases}
        A(x)\geq0\\
        B(x)\geq0\\
        A(x)\leq B^n(x)
    \end{cases}
\end{equation}
Esempio: $\sqrt{x}<x-2$

$\begin{array}{l}
   \begin{cases}
       x\geq0\\
       x-2>0\\
       x<(x-2)^2
   \end{cases}\implies\begin{cases}
       x\geq0\\
       x>2\\
       x<1\lor x>4
   \end{cases}
\end{array}$

\begin{tikzpicture}[x=0.5cm,y=0.5cm,font=\footnotesize,
  open/.style={draw,fill=white,circle,inner sep=0pt,minimum size=4pt},
  closed/.style={draw,fill=black,circle,inner sep=0pt,minimum size=4pt}]

  \def\xa{3}
  \def\xb{7}
  \def\xc{11}
  \def\xd{15}
  \def\xe{19}

  \fill[gray!20] (\xd,0) rectangle (\xe,-7);

  \draw[-{Stealth[length=2mm]}, thick] (-0.5,0) -- (20,0) node[below left] {$x$};

  \node[above] at (\xa,0.15) {$0$};
  \node[above] at (\xb,0.15) {$1$};
  \node[above] at (\xc,0.15) {$2$};
  \node[above] at (\xd,0.15) {$4$};

  \foreach \x in {\xa,\xb,\xc,\xd}
    \draw[very thick] (\x,0.2) -- (\x,-7);

  \draw[thick,dashed] (-0.5,-2) -- (\xa,-2);
  \draw[thick]        (\xa,-2) -- (\xe,-2);
  \node[closed] at (\xa,-2) {};

  \draw[thick,dashed] (-0.5,-4) -- (\xc,-4);
  \draw[thick]        (\xc,-4) -- (\xe,-4);
  \node[open] at (\xc,-4) {};

  \draw[thick]        (-0.5,-6) -- (\xb,-6);
  \draw[thick,dashed] (\xb,-6) -- (\xd,-6);
  \draw[thick]        (\xd,-6) -- (\xe,-6);
  \node[open] at (\xb,-6) {};
  \node[open] at (\xd,-6) {};

\end{tikzpicture}$\implies x>4$

Se il secondo membro è un numero $k\in\R$ allora esistono delle tecniche per sveltire i calcoli:
\begin{enumerate}
    \item $\sqrt{A(x)}<k$, con $k<0$: $\nexists x\in\R$
    \item $\sqrt{A(x)}\leq k$, con $k=0$: $\sqrt{A(x)}=0$
    \item $\sqrt{A(x)}<k$, con $k>0$: $\begin{cases}
        A(x)\geq0\\
        A(x)<k^2
    \end{cases}$
\end{enumerate}
Per risolvere disequazioni irrazionali nella forma $\sqrt[n]{A(x)}>B(x)$ con $n$ pari si riconducono al seguente metodo:
\begin{equation}
    \sqrt[n]{A(x)}>B(x)\iff\begin{cases}
        B(x)<0\\
        A(x)\geq0
    \end{cases}\cup\begin{cases}
        B(x)\geq0\\
        A(x)>B^n(x)
    \end{cases}
\end{equation}
In caso venisse presa in considerazione anche l'uguaglianza si ricorre al seguente metodo:
\begin{equation}
    \sqrt[n]{A(x)}\geq B(x)\iff\begin{cases}
        B(x)<0\\
        A(x)\geq0
    \end{cases}\cup\begin{cases}
        B(x)\geq0\\
        A(x)\geq B^n(x)
    \end{cases}
\end{equation}
Esempio: $\sqrt{9-x^2}>x+3$

$\begin{array}{l}
   \begin{cases}
        x+3<0\\
        9-x^2\geq0
   \end{cases}\cup\begin{cases}
       x+3\geq0\\
       9-x^2>(x+3)^2
   \end{cases}\\
   \begin{cases}
       x<-3\\
       -3\leq x\leq3
   \end{cases}\cup\begin{cases}
       x\geq-3\\
       -3<x<0
   \end{cases}
\end{array}$

\begin{tikzpicture}[x=0.5cm,y=0.5cm,font=\footnotesize,
  open/.style={draw,fill=white,circle,inner sep=0pt,minimum size=4pt},
  closed/.style={draw,fill=black,circle,inner sep=0pt,minimum size=4pt}]

  \begin{scope}
    \def\xa{3}
    \def\xb{7}

    \draw[-{Stealth[length=2mm]}, thick] (-0.5,0) -- (9.5,0) node[below left] {$x$};

    \node[above] at (\xa,0.15) {$-3$};
    \node[above] at (\xb,0.15) {$3$};

    \foreach \x in {\xa,\xb}
      \draw[very thick] (\x,0.2) -- (\x,-5);

    \draw[thick,dashed] (-0.5,-2) -- (\xa,-2);
    \draw[thick]        (\xa,-2) -- (\xb,-2);
    \draw[thick,dashed] (\xb,-2) -- (9,-2);
    \node[closed] at (\xa,-2) {};
    \node[closed] at (\xb,-2) {};

    \draw[thick]        (-0.5,-4) -- (\xa,-4);
    \draw[thick,dashed] (\xa,-4) -- (9,-4);
    \node[open] at (\xa,-4) {};

    \node at (4.25,-6.5) {Primo sistema};
  \end{scope}

  \begin{scope}[xshift=7cm]
    \def\xa{3}
    \def\xb{7}

    \fill[gray!20] (\xa,0) rectangle (\xb,-5);

    \draw[-{Stealth[length=2mm]}, thick] (-0.5,0) -- (9.5,0) node[below left] {$x$};

    \node[above] at (\xa,0.15) {$-3$};
    \node[above] at (\xb,0.15) {$0$};

    \foreach \x in {\xa,\xb}
      \draw[very thick] (\x,0.2) -- (\x,-5);

    \draw[thick,dashed] (-0.5,-2) -- (\xa,-2);
    \draw[thick]        (\xa,-2) -- (9,-2);
    \node[closed] at (\xa,-2) {};

    \draw[thick,dashed] (-0.5,-4) -- (\xa,-4);
    \draw[thick]        (\xa,-4) -- (\xb,-4);
    \draw[thick,dashed] (\xb,-4) -- (9,-4);
    \node[open] at (\xa,-4) {};
    \node[open] at (\xb,-4) {};

    \node at (4.25,-6.5) {Secondo sistema};
  \end{scope}

\end{tikzpicture}

$\implies-3<x<0$

Se il secondo membro è un numero $k\in\R$ allora esistono delle tecniche per sveltire i calcoli:
\begin{enumerate}
    \item $\sqrt{A(x)}>k$, con $k<0$: $A(x)\geq0$
    \item $\sqrt{A(x)}>k$, con $k=0$: $A(x)>0$
    \item $\sqrt{A(x)}>k$, con $k>0$: $A(x)>k^2$
\end{enumerate}
In caso compaia più di un radicale occorre procedere nello stesso modo delle equazioni.
\newpage
\section{Equazioni e disequazioni con valori assoluti}
\begin{definizione}
    Dato $x\in \R$, si chiama valore assoluto di $x$, e si indica con il simbolo $|x|$, il numero reale così definito:
    \begin{equation}
        |x|=\begin{cases}
            x&\text{se }x\geq0\\
            -x&\text{se }x<0
        \end{cases}
    \end{equation}
\end{definizione}
\begin{definizione}
    Data un'espressione algebrica $f(x)$, si chiama valore assoluto di $f(x)$, e si indica con il simbolo $|f(x)|$, l'espressione così definita:
    \begin{equation}
        |f(x)|=\begin{cases}
            f(x)&\text{se }f(x)\geq0\\
            -f(x)&\text{se }f(x)<0
        \end{cases}
    \end{equation}
\end{definizione}
Di seguito sono elencate le principali proprietà del valore assoluto:
\begin{enumerate}
    \item $|-x|=|x|\quad \forall x\in\R$
    \item $|x|^2=x^2\quad \forall x\in\R$
    \item $\sqrt{x^2}=|x|\quad \forall x\in\R$
    \item $|x+y|\leq|x|+|y|\quad\forall x,y\in\R$
    \item $|xy|=|x||y|\quad\forall x,y\in\R$
    \item $\left|\dfrac{x}{y}\right|=\dfrac{|x|}{|y|}\quad\forall x,y\in\R\land y\ne0$
    \item $|x|=|y|\iff x=\pm y\quad\forall x,y\in\R$
\end{enumerate}
Per risolvere equazioni del tipo $|A(x)|=B(x)$ si ricorre al seguente metodo:
\begin{equation}
    |A(x)|=B(x)\iff\begin{cases}
        A(x)\geq0\\
        A(x)=B(x)
    \end{cases}\cup\begin{cases}
        A(x)<0\\
        -A(x)=B(x)
    \end{cases}
\end{equation}
Se il secondo membro è un numero $k\in\R$ allora esistono delle tecniche per sveltire i calcoli:
\begin{enumerate}
    \item $|A(x)|=k$, con $k<0$: $\nexists x\in\R$
    \item $|A(x)|=k$, con $k=0$: $A(x)=0$
    \item $|A(x)|=k$, con $k>0$: $A(x)=\pm k$
\end{enumerate}
Per risolvere equazioni del tipo $|A(x)|=|B(x)|$ si sfrutta la $7^a$ proprietà dei moduli:
\begin{equation}
    |A(x)|=|B(x)|\iff A(x)=B(x)\lor A(x)=-B(x)
\end{equation}
Se invece compare più di un modulo in un'equazione bisogna studiare il segno di ogni modulo e costruire tanti sistemi quanti sono gli intervalli trovati dallo studio.

Per risolvere disequazioni del tipo $|A(x)|<B(x)$ si ricorre al seguente metodo:
\begin{equation}
    |A(x)|<B(x)\iff \begin{cases}
        A(x)<B(x)\\
        A(x)>-B(x)
    \end{cases}
\end{equation}
Per risolvere disequazioni del tipo $|A(x)|>B(x)$ si ricorre al seguente metodo:
\begin{equation}
    |A(x)|>B(x)\iff A(x)<-B(x)\lor A(x)>B(x)
\end{equation}
Per risolvere disequazioni del tipo $|A(x)|\lessgtr|B(x)|$ si sfrutta la $2^a$ proprietà dei moduli:
\begin{equation}
    |A(x)|\lessgtr|B(x)|\iff A^2(x)\lessgtr B^2(x)
\end{equation}
Se invece compare più di un modulo bisogna ricondursi allo stesso metodo delle equazioni con più moduli.
\newpage
\section{Funzioni}
\subsection{Introduzione alle funzioni}
\begin{definizione}
    Siano $X$ e $Y$ due insiemi; si dice funzione da $X$ a $Y$ una relazione che associa a ogni elemento di $X$ un solo elemento di $Y$.
\end{definizione}
L'insieme di partenza $X$ si chiama dominio, mentre l'insieme di arrivo $Y$ si chiama codominio. Una funzione $f$ con dominio $X$ e codominio $Y$ si definisce tramite la scrittura $f:X\to Y$. L'immagine di un elemento $x$, l'elemento nel codominio che tramite $f$ corrisponde a $x$, viene indicata con $f(x)$. Simmetricamente si può dire che $x$ è la controimmagine di $y$, se $y$ è l'immagine di $x$. L'insieme immagine di una funzione viene indicato con $f(X)$.

Le funzioni reali di variabile reale si possono indicare tramite un'equazione del tipo:
\begin{equation}
    y=f(x)
\end{equation}
dove $f(x)$ è l'espressione analitica della funzione, oppure mediante la seguente scrittura:
\begin{equation}
    f(x)=\dots
\end{equation}
dove al posto dei puntini c'è l'espressione analitica della funzione.

Esempio 1: $f:\R\to\R_0^+$, funzione che associa a ogni numero reale il suo quadrato, può essere definita in una delle seguenti forme:

$y=x^2\quad f(x)=x^2$

Esempio 2: $f:\R\to\R$, funzione che associa a ogni numero reale sé stesso, detta funzione identità, può essere definita in una delle seguenti forme:

$y=x\quad f(x)=x$

In un'equazione $y=f(x)$, $x$ è la variabile indipendente e $y$ la variabile dipendente da $x$.

Di seguito è riportata la suddivisione delle funzioni reali di variabile reale:
\begin{enumerate}
    \item Algebriche
    \begin{enumerate}
        \item Intere
        \begin{enumerate}
            \item Razionali, es: $y=x^4-x^2$
            \item Irrazionali, es: $y=\sqrt{x^2+3x}$
        \end{enumerate}
        \item Frazionarie
        \begin{enumerate}
            \item Razionali, es: $y=\dfrac{1}{x^4-x^2}$
            \item Irrazionali, es: $y=\dfrac{1}{\sqrt[3]{x+1}}$
        \end{enumerate}
    \end{enumerate}
    \item Trascendenti, es: $y=x\cdot 6^x$
\end{enumerate}
Per calcolare il dominio di una funzione reale di variabile reale bisogna porre ogni denominatore diverso da $0$ e ogni radicando, di radicale con indice pari, maggiore o uguale a $0$.

Esempio 1: $y=x^4-x^2$

Non ci sono né denominatori né radicali con indici pari, quindi il dominio della funzione è $\R$.

Esempio 2: $y=\dfrac{\sqrt{x-2}}{x^2+x+1}$

Bisogna porre $x-2\geq0$ e $x^2+x+1\ne0$ contemporaneamente ricorrendo quindi a un sistema:

$\begin{cases}
    x-2\geq0\\
    x^2+x+1\ne0
\end{cases}\implies x\geq2$

Quindi il dominio della funzione è l'intervallo $[2,+\infty)$.

Il grafico di una funzione $f:X\to Y$ è l'insieme $\set{(x,f(x)):x\in X}$.

Per capire se un grafico è un grafico di una funzione bisogna tracciare delle rette parallele all'asse $y$. Se almeno una retta interseca il grafico in più di un punto allora il grafico non rappresenta una funzione.
\newpage
Esempio 1: Grafico di una funzione

\begin{center}
\begin{tikzpicture}[scale=0.6, >=Stealth]

  \draw[->] (-2.2,0) -- (4.6,0) node[below] {$x$};
  \draw[->] (0,-2.8) -- (0,3.6) node[left] {$y$};
  \node[above left] at (0,0) {$O$};

  \begin{scope}
    \clip (-2.1,-2.7) rectangle (4.5,3.5);
    \draw[thick, smooth, samples=100, domain=-1.5:4.5]
      plot (\x, {0.5*\x*\x*(\x-3)});
  \end{scope}

  \foreach \a in {-1,1,2}{
    \draw[dashed] (\a,-2.8) -- (\a,3.6);
  }

  \foreach \a in {-1,1,2}{
    \fill (\a, {0.5*\a*\a*(\a-3)}) circle (2pt);
  }

\end{tikzpicture}
\end{center}
Esempio 2: Grafico che non rappresenta una funzione

\begin{center}
\begin{tikzpicture}[scale=0.6, >=Stealth]

  \draw[->] (-4.2,0) -- (4.6,0) node[below] {$x$};
  \draw[->] (0,-4.2) -- (0,4.4) node[left] {$y$};
  \node[below left, inner sep=2pt] at (0,0) {$O$};

  \begin{scope}
    \clip (-4.1,-4.1) rectangle (4.1,4.1);
    \draw[thick, smooth, samples=100, domain=-4.2:4.2]
      plot (\x, {sqrt(1+\x*\x)});
    \draw[thick, smooth, samples=100, domain=-4.2:4.2]
      plot (\x, {-sqrt(1+\x*\x)});
  \end{scope}

  \draw[dashed] (-2,-4.2) -- (-2,4.2);

  \fill (-2, {sqrt(5)})  circle (2pt);
  \fill (-2, {-sqrt(5)}) circle (2pt);

\end{tikzpicture}
\end{center}
Per tracciare un grafico sapendo l'espressione analitica della funzione si può procedere a punti sostituendo alla $x$ dei valori opportuni alla costruzione del grafico.

\begin{definizione}
    Due funzioni $f$ e $g$ sono uguali se hanno lo stesso dominio $X$ e inoltre risulta:
    \begin{equation}
        f(x)=g(x)\quad\forall x\in X
    \end{equation}
\end{definizione}
Esempio 1: $f(x)=x^2-1$ e $g(x)=\dfrac{x^4-1}{x^2+1}$

Il dominio è uguale tra le due funzioni, il dominio è $\R$.

Se poniamo $f(x)=g(x)$:

$\begin{array}{l}
    x^2-1=\dfrac{x^4-1}{x^2+1}\\
    x^2-1=\dfrac{(x^2+1)(x^2-1)}{x^2+1}\\
    x^2-1=x^2-1\\
    \forall x\in\R
\end{array}$

Le due funzioni sono uguali per ogni $x$ nel dominio, quindi $f=g$.

Esempio 2: $f(x)=x+1$ e $g(x)=\dfrac{x^2-1}{x-1}$

Il dominio di $f(x)$ è $\R$, mentre il dominio di $g(x)$ è $\R\setminus\set{1}$ quindi $f\ne g$.
\newpage
\subsection{Prime proprietà delle funzioni reali di variabile reale}
Studiare il segno di una funzione $f$ significa stabilire per quali valori del dominio $f(x)>0$, $f(x)=0$ e $f(x)<0$. Per i valori di $x$ tali che $f(x)>0$ la funzione si trova al di sopra dell'asse delle ascisse, per $f(x)=0$ si trovano gli zeri della funzione e per $f(x)<0$ la funzione si trova sotto l'asse delle ascisse.
\begin{center}
\begin{tikzpicture}[scale=0.55]

    \fill[pattern=north east lines, pattern color=gray!65] (-4.5,0) rectangle (-4,3.5);
    \fill[pattern=north east lines, pattern color=gray!65] (-4,-2.5) rectangle (1,0);
    \fill[pattern=north east lines, pattern color=gray!65] (1,0) rectangle (3,3.5);
    \fill[pattern=north east lines, pattern color=gray!65] (3,-2.5) rectangle (4.2,0);

    \draw[dashed, thick, black!80] (-4,-2.5) -- (-4,3.5);
    \draw[dashed, thick, black!80] (1,-2.5) -- (1,3.5);
    \draw[dashed, thick, black!80] (3,-2.5) -- (3,3.5);

    \draw[-{Stealth[length=2.5mm]}, thick] (-4.8,0) -- (4.8,0) node[below left, font=\small] {$x$};
    \draw[-{Stealth[length=2.5mm]}, thick] (0,-2.8) -- (0,4.0) node[below left, font=\small] {$y$};

    \draw[very thick, black, smooth, samples=100, domain=-4.1:3.15]
        plot (\x, {0.1*(\x+4)*(\x-1)*(\x-3)});

    \draw[very thick, dashed, black, smooth, samples=15, domain=-4.2:-4.1]
        plot (\x, {0.1*(\x+4)*(\x-1)*(\x-3)});
    \draw[very thick, dashed, black, smooth, samples=15, domain=3.15:3.25]
        plot (\x, {0.1*(\x+4)*(\x-1)*(\x-3)});

    \fill (-4, 0) circle (3pt);
    \fill ( 1, 0) circle (3pt);
    \fill ( 3, 0) circle (3pt);

    \node[below left, fill=white, inner sep=1pt, font=\small] at (-4,-0.1) {$-4$};
    \node[below left, fill=white, inner sep=1pt, font=\small] at (0,-0.1) {$O$};
    \node[below right, fill=white, inner sep=1pt, font=\small] at (1,-0.1) {$1$};
    \node[below right, fill=white, inner sep=1pt, font=\small] at (3,-0.1) {$3$};

    \node[above left, font=\small] at (2.5, 2.2) {$y = f(x)$};

\end{tikzpicture}
\end{center}

Le funzioni inoltre possono essere pari o dispari o nessuna delle due.
\begin{definizione}
    Una funzione definita da $y=f(x)$ si dice pari se per ogni $x$ appartenente al dominio della funzione anche $-x$ vi appartiene e risulta:
    \begin{equation}
        f(-x)=f(x)
    \end{equation}
    In tal caso, il grafico della funzione risulta simmetrico rispetto all'asse $y$.
\end{definizione}
\begin{definizione}
    Una funzione definita da $y=f(x)$ si dice dispari se per ogni $x$ appartenente al dominio della funzione anche $-x$ vi appartiene e risulta:
    \begin{equation}
        f(-x)=-f(x)
    \end{equation}
    In tal caso, il grafico della funzione risulta simmetrico rispetto all'origine.
\end{definizione}
\begin{center}
\begin{tikzpicture}[scale=0.55]

    \begin{scope}
        \draw[-{Stealth[length=2.5mm]}, thick] (-3.2,0) -- (3.2,0) node[below left, font=\small] {$x$};
        \draw[-{Stealth[length=2.5mm]}, thick] (0,-2.2) -- (0,4.5) node[below left, font=\small] {$y$};
        \node[above left, font=\small] at (0,0) {$O$};

        \draw[very thick, black, smooth, samples=100, domain=-1.9:1.9]
            plot (\x, {0.35*(\x)^4 - 1.4*(\x)^2});
        \draw[very thick, dashed, black, smooth, samples=15, domain=-2.05:-1.9]
            plot (\x, {0.35*(\x)^4 - 1.4*(\x)^2});
        \draw[very thick, dashed, black, smooth, samples=15, domain=1.9:2.05]
            plot (\x, {0.35*(\x)^4 - 1.4*(\x)^2});

        \coordinate (P) at (1.6, {0.35*(1.6)^4 - 1.4*(1.6)^2});
        \coordinate (P1) at (-1.6, {0.35*(-1.6)^4 - 1.4*(-1.6)^2});

        \draw[dashed, thick] (P1) -- (P);
        \fill (P) circle (3pt) node[right, font=\small] {$P$};
        \fill (P1) circle (3pt) node[left, font=\small] {$P'$};

        \node[font=\bfseries] at (0, -3.2) {funzione pari};
    \end{scope}

    \begin{scope}[xshift=9cm]
        \draw[-{Stealth[length=2.5mm]}, thick] (-3.2,0) -- (3.2,0) node[below left, font=\small] {$x$};
        \draw[-{Stealth[length=2.5mm]}, thick] (0,-3.2) -- (0,3.2) node[below left, font=\small] {$y$};
        \node[above left, font=\small] at (0,0) {$O$};

        \draw[very thick, black, smooth, samples=100, domain=-2.1:2.1]
            plot (\x, {0.35*(\x)^3 - 1.5*(\x)});
        \draw[very thick, dashed, black, smooth, samples=15, domain=-2.25:-2.1]
            plot (\x, {0.35*(\x)^3 - 1.5*(\x)});
        \draw[very thick, dashed, black, smooth, samples=15, domain=2.1:2.25]
            plot (\x, {0.35*(\x)^3 - 1.5*(\x)});

        \coordinate (Q) at (1.7, {0.35*(1.7)^3 - 1.5*(1.7)});
        \coordinate (Q1) at (-1.7, {0.35*(-1.7)^3 - 1.5*(-1.7)});

        \draw[dashed, thick] (Q1) -- (Q);
        \fill (Q) circle (3pt) node[right, font=\small] {$P$};
        \fill (Q1) circle (3pt) node[left, font=\small] {$P'$};

        \node[font=\bfseries] at (0, -4.2) {funzione dispari};
    \end{scope}

\end{tikzpicture}
\end{center}
\newpage
Inoltre si può anche studiare la monotonia di una funzione.
\begin{center}
\begin{tikzpicture}[scale=0.55]

    \draw[-{Stealth[length=2.5mm]}, thick] (-8.5,0) -- (11.5,0) node[below left, font=\small] {$x$};
    \draw[-{Stealth[length=2.5mm]}, thick] (0,-2.5) -- (0,7.5) node[below left, font=\small] {$y$};
    \node[above left, font=\small] at (0,0) {$O$};

    \draw[dashed, thin, black!60] (-7,0) -- (-7, 1);
    \draw[dashed, thin, black!60] (-4,0) -- (-4, 5);
    \draw[dashed, thin, black!60] (2,0)  -- (2, -1);
    \draw[dashed, thin, black!60] (5,0)  -- (5, 6);
    \draw[dashed, thin, black!60] (10,0) -- (10, 6);

    \draw[very thick, black, ->-={0.5}] (-7, 1) -- (-4, 5);

    \draw[very thick, black, smooth, samples=100, domain=-4:2]
        plot (\x, { (1/18)*(\x)^3 + (1/6)*(\x)^2 - (4/3)*(\x) + (5/9) });
    \draw[very thick, black, -{Stealth[length=2.5mm]}, samples=5, domain=-1.1:-1]
        plot (\x, { (1/18)*(\x)^3 + (1/6)*(\x)^2 - (4/3)*(\x) + (5/9) });

    \draw[very thick, black, smooth, samples=100, domain=2:5]
        plot (\x, { (-14/27)*(\x-2)^3 + (7/3)*(\x-2)^2 - 1 });
    \draw[very thick, black, -{Stealth[length=2.5mm]}, samples=5, domain=3.4:3.5]
        plot (\x, { (-14/27)*(\x-2)^3 + (7/3)*(\x-2)^2 - 1 });

    \draw[very thick, black, ->-={0.5}] (5, 6) -- (10, 6);

    \fill (-7, 1)  circle (3pt) node[above left, font=\small] {$A$};
    \fill (-4, 5)  circle (3pt) node[above, font=\small] {$B$};
    \fill ( 2, -1) circle (3pt) node[below left, font=\small] {$C$};
    \fill ( 5, 6)  circle (3pt) node[above left, font=\small] {$D$};
    \fill (10, 6)  circle (3pt) node[above, font=\small] {$E$};

    \node[below left, fill=white, inner sep=1pt, font=\small] at (-7,-0.1) {$-7$};
    \node[below left, fill=white, inner sep=1pt, font=\small] at (-4,-0.1) {$-4$};
    \node[above left, fill=white, inner sep=1pt, font=\small] at (2,0) {$2$};
    \node[below right, fill=white, inner sep=1pt, font=\small] at (5,-0.1) {$5$};
    \node[below right, fill=white, inner sep=1pt, font=\small] at (10,-0.1) {$10$};

\end{tikzpicture}
\end{center}
Questa funzione è strettamente crescente in $[A,B]$, strettamente decrescente in $[B,C]$, strettamente crescente in $[C,D]$ e costante in $[D,E]$.
\begin{definizione}
    Sia $I\subseteq D(f)$.
    \begin{enumerate}
        \item $f$ si dice strettamente crescente in $I$ se:
        \begin{equation}
            x_1<x_2\implies f(x_1)<f(x_2)\quad \forall x_1,x_2\in I
        \end{equation}
        \item $f$ si dice strettamente decrescente in $I$ se:
        \begin{equation}
            x_1<x_2\implies f(x_1)>f(x_2)\quad \forall x_1,x_2\in I
        \end{equation}
        \item $f$ si dice costante in $I$ se:
        \begin{equation}
            f(x_1)=f(x_2)\quad \forall x_1,x_2\in I
        \end{equation}
    \end{enumerate}
\end{definizione}
\begin{definizione}
    Sia $I\subseteq D(f)$.
    \begin{enumerate}
        \item $f$ si dice crescente in senso lato in $I$ se:
        \begin{equation}
            x_1<x_2\implies f(x_1)\leq f(x_2)\quad \forall x_1,x_2\in I
        \end{equation}
        \item $f$ si dice decrescente in senso lato in $I$ se:
        \begin{equation}
            x_1<x_2\implies f(x_1)\geq f(x_2)\quad \forall x_1,x_2\in I
        \end{equation}
    \end{enumerate}
\end{definizione}
\begin{definizione}
    Una funzione crescente o decrescente in un insieme $I$ viene detta monotòna in $I$.
\end{definizione}
\newpage
\subsection{Funzioni iniettive, suriettive, biiettive}
Le funzioni si possono classificare in base al loro comportamento rispetto al codominio.
\begin{definizione}
    Una funzione $f:X\to Y$ si dice iniettiva se ogni elemento di $Y$ ha al massimo una controimmagine in $X$.
\end{definizione}
\begin{figure}[htbp]
  \centering

  \begin{minipage}[t]{0.45\textwidth}
    \centering
    \begin{tikzpicture}[scale=0.45, >=Latex]

      \draw[thick, ->] (-6.5, 0) -- (2.2, 0) node[below left] {$x$};
      \draw[thick, ->] (0, -2.5) -- (0, 5.5) node[below left] {$y$};
      \node[below left] at (0, 0) {$O$};

      \draw[dashed, thick] (-6.5, -1.5) -- (1.8, -1.5);
      \draw[dashed, thick] (-6.5, 2.0) -- (1.8, 2.0);
      \draw[dashed, thick] (-6.5, 3.5) -- (1.8, 3.5);

      \draw[very thick, domain=-5.5:1.6, samples=100] plot (\x, {2.5 * (1.4^\x)});

      \draw[very thick, dashed, domain=-6.3:-5.5, samples=20] plot (\x, {2.5 * (1.4^\x)});
      \draw[very thick, dashed, domain=1.6:1.95, samples=20] plot (\x, {2.5 * (1.4^\x)});

      \node at (-3.8, 1.3) {$y = f(x)$};
    \end{tikzpicture}
    \caption{La funzione $f(x)$ è iniettiva.}
  \end{minipage}
  \hfill

  \begin{minipage}[t]{0.45\textwidth}
    \centering
    \begin{tikzpicture}[scale=0.45, >=Latex]

      \draw[thick, ->] (-4.5, 0) -- (4.5, 0) node[below left] {$x$};
      \draw[thick, ->] (0, -3.8) -- (0, 4.5) node[below left] {$y$};
      \node[below left] at (0, 0) {$O$};

      \draw[dashed, thick] (-4.3, 2.0) -- (4.3, 2.0);
      \draw[dashed, thick] (-4.3, -2.0) -- (4.3, -2.0);
      \draw[dashed, thick] (-4.3, -3.0) -- (4.3, -3.0);

      \draw[very thick, smooth, domain=-2.1:2.1, samples=150] plot (\x, {0.8 * \x * \x});

      \draw[very thick, dashed, smooth, domain=-2.3:-2.1, samples=20] plot (\x, {0.8 * \x * \x});
      \draw[very thick, dashed, smooth, domain=2.1:2.3, samples=20] plot (\x, {0.8 * \x * \x});

      \node at (2.8, 1.0) {$y = g(x)$};
    \end{tikzpicture}
    \caption{La funzione $g(x)$ non è iniettiva.}
  \end{minipage}
\end{figure}
\begin{definizione}
    Una funzione $f:X\to Y$ si dice suriettiva se ogni elemento di $Y$ ha almeno una controimmagine in $X$.
\end{definizione}
\begin{figure}[htbp]
  \centering

  \begin{minipage}[t]{0.45\textwidth}
    \centering
    \begin{tikzpicture}[scale=0.45, >=Latex]

      \draw[thick, ->] (-4.5, 0) -- (4.5, 0) node[below left] {$x$};
      \draw[thick, ->] (0, -4.5) -- (0, 4.5) node[below left] {$y$};
      \node[above left] at (0, 0) {$O$};

      \draw[dashed, thick] (-4.3, 1.5) -- (4.3, 1.5);
      \draw[dashed, thick] (-4.3, -0.5) -- (4.3, -0.5);
      \draw[dashed, thick] (-4.3, -3.5) -- (4.3, -3.5);

      \draw[very thick, smooth, domain=-1.95:3.8, samples=150] plot (\x, {0.25 * \x * \x * \x - 0.75 * \x * \x});

      \draw[very thick, dashed, smooth, domain=-2.15:-1.95, samples=20] plot (\x, {0.25 * \x * \x * \x - 0.75 * \x * \x});
      \draw[very thick, dashed, smooth, domain=3.8:4.0, samples=20] plot (\x, {0.25 * \x * \x * \x - 0.75 * \x * \x});

      \fill (3.47, 1.5) circle (3.5pt);
      \fill (-0.73, -0.5) circle (3.5pt);
      \fill (1.0, -0.5) circle (3.5pt);
      \fill (2.73, -0.5) circle (3.5pt);
      \fill (-1.78, -3.5) circle (3.5pt);

      \node at (-3.1, -1.8) {$y = f(x)$};
    \end{tikzpicture}
    \caption{La funzione $f(x)$ è suriettiva.}
  \end{minipage}
  \hfill

  \begin{minipage}[t]{0.45\textwidth}
    \centering
    \begin{tikzpicture}[scale=0.45, >=Latex]

      \draw[thick, ->] (-4.5, 0) -- (4.5, 0) node[below left] {$x$};
      \draw[thick, ->] (0, -3.8) -- (0, 4.5) node[below left] {$y$};
      \node[below left] at (0, 0) {$O$};

      \draw[dashed, thick] (-4.3, 2.0) -- (4.3, 2.0);
      \draw[dashed, thick] (-4.3, -2.0) -- (4.3, -2.0);
      \draw[dashed, thick] (-4.3, -3.0) -- (4.3, -3.0);

      \draw[very thick, smooth, domain=-2.1:2.1, samples=150] plot (\x, {0.8 * \x * \x});

      \draw[very thick, dashed, smooth, domain=-2.3:-2.1, samples=20] plot (\x, {0.8 * \x * \x});
      \draw[very thick, dashed, smooth, domain=2.1:2.3, samples=20] plot (\x, {0.8 * \x * \x});

      \fill (-1.58, 2.0) circle (3.5pt);
      \fill (1.58, 2.0) circle (3.5pt);

      \node at (2.8, 1.0) {$y = g(x)$};
    \end{tikzpicture}
    \caption{La funzione $g(x)$ non è suriettiva.}
  \end{minipage}
\end{figure}
\begin{definizione}
    Una funzione $f:X\to Y$ che è sia iniettiva sia suriettiva si dice biiettiva.
\end{definizione}
\begin{figure}[htbp]
  \centering

  \begin{minipage}[t]{0.45\textwidth}
    \centering
    \begin{tikzpicture}[scale=0.45, >=Latex]

      \draw[thick, ->] (-4.5, 0) -- (4.5, 0) node[below left] {$x$};
      \draw[thick, ->] (0, -4.5) -- (0, 4.5) node[below left] {$y$};
      \node[below left] at (0, 0) {$O$};

      \draw[dashed, thick] (-4.3, 3.0) -- (4.3, 3.0);
      \draw[dashed, thick] (-4.3, 1.8) -- (4.3, 1.8);
      \draw[dashed, thick] (-4.3, -2.5) -- (4.3, -2.5);

      \draw[very thick, smooth, domain=-3.0:3.0, samples=150] plot (\x, {0.15 * \x * \x * \x});

      \draw[very thick, dashed, smooth, domain=-3.15:-3.0, samples=20] plot (\x, {0.15 * \x * \x * \x});
      \draw[very thick, dashed, smooth, domain=3.0:3.15, samples=20] plot (\x, {0.15 * \x * \x * \x});

      \node at (2.5, -1.2) {$y = f(x)$};
    \end{tikzpicture}
    \caption{La funzione $f(x)$ è biiettiva.}
  \end{minipage}
  \hfill

  \begin{minipage}[t]{0.45\textwidth}
    \centering
    \begin{tikzpicture}[scale=0.45, >=Latex]

      \draw[thick, ->] (-4.5, 0) -- (4.5, 0) node[below left] {$x$};
      \draw[thick, ->] (0, -4.5) -- (0, 4.5) node[below left] {$y$};
      \node[below left] at (0, 0) {$O$};

      \draw[dashed, thick] (-4.3, 1.8) -- (4.3, 1.8);
      \draw[dashed, thick] (-4.3, -1.8) -- (4.3, -1.8);

      \draw[very thick, smooth, domain=-2.25:2.25, samples=150] plot (\x, {0.7 * \x * \x * \x - 2.6 * \x});

      \draw[very thick, dashed, smooth, domain=-2.35:-2.25, samples=20] plot (\x, {0.7 * \x * \x * \x - 2.6 * \x});
      \draw[very thick, dashed, smooth, domain=2.25:2.35, samples=20] plot (\x, {0.7 * \x * \x * \x - 2.6 * \x});

      \node at (2.8, -3.1) {$y = g(x)$};
    \end{tikzpicture}
    \caption{La funzione $g(x)$ non è biiettiva.}
  \end{minipage}
\end{figure}
\newpage
Questa classificazione può essere vista anche mediante la seguente rappresentazione insiemistica.
\begin{figure}[htbp]
  \centering
  \begin{tikzpicture}[scale=0.85]

    \draw[thick] (-4.5, -2.5) rectangle (4.5, 2.5);
    \node at (0, 2.8) {\textbf{Funzioni}};

    \draw[thick] (-1.2, 0) ellipse (2.4cm and 1.5cm);
    \draw[thick] (1.2, 0) ellipse (2.4cm and 1.5cm);

    \node at (-2.3, 0) {iniettive};
    \node at (0, 0) {biiettive};
    \node at (2.3, 0) {suriettive};
    \node at (0, -1.9) {né iniettive né suriettive};
  \end{tikzpicture}
  \caption{Diagramma di Venn delle funzioni.}
\end{figure}
\newpage
\subsection{Funzioni composte}
\begin{definizione}
    Date due funzioni $f$ e $g$, si dice funzione composta di $f$ e $g$, e si indica con il simbolo $g\circ f$, la funzione definita da:
    \begin{equation}
        (g\circ f)(x)=g(f(x))
    \end{equation}
\end{definizione}
Il dominio di $g\circ f$ è costituito da tutti gli elementi appartenenti al dominio di $f$ tali che $f(x)$ appartiene al dominio di $g$.

ESEMPIO: $f(x)=x-1\quad g(x)=\sqrt{x+3}$

$\begin{array}{ll}
    (g\circ f)(x)=g(f(x))=g(x-1)=\sqrt{(x-1)+3}=\sqrt{x+2} & D=[-2,+\infty[ \\
    (f\circ g)(x)=f(g(x))=f(\sqrt{x+3})=\sqrt{x+3}-1 & D=[-3,+\infty[
\end{array}$

Ciò mostra che $f\circ g\ne g\circ f$, quindi la composizione di due funzioni non è commutativa.
\newpage
\subsection{Funzione inversa}

\noindent
\begin{minipage}[t]{0.32\textwidth}
\vspace{0pt}
\centering
\begin{tikzpicture}[scale=0.78, >=Stealth, thick]

    \draw[->] (-0.3,0) -- (4.1,0) node[right] {$x$};
    \draw[->] (0,-0.3) -- (0,3.5) node[above] {$y$};
    \node[below left] at (0,0) {$O$};

    \node at (2,4.5) {$f:[a,b]\longrightarrow\mathbb{R}$};

    \draw plot[smooth] coordinates {
        (0.6,3) (1.1,2.9) (1.6,2.6)
        (2.1,2.0) (2.6,1.2) (3.2,0.5)
    };

    \fill (0.6,3) circle (2pt);
    \fill (3.2,0.5) circle (2pt);
    \fill (2.1,2.0) circle (2pt);

    \draw[dashed] (0.6,3) -- (0.6,0);
    \draw[dashed] (0,3) -- (0.6,3);
    \draw[dashed] (3.2,0.5) -- (3.2,0);
    \draw[dashed] (0,0.5) -- (3.2,0.5);

    \draw[->] (2.1,0) -- (2.1,2.0);
    \draw[->] (2.1,2.0) -- (0,2.0);

    \node[left] at (0,3) {$d$};
    \node[left] at (0,2.0) {$y_0$};
    \node[left] at (0,0.5) {$c$};

    \node[below] at (0.6,0) {$a$};
    \node[below] at (2.1,0) {$x_0$};
    \node[below] at (3.2,0) {$b$};

\end{tikzpicture}
\end{minipage}
\hfill
\begin{minipage}[t]{0.32\textwidth}
\vspace{0pt}
\centering
\begin{tikzpicture}[scale=0.78, >=Stealth, thick]

    \draw[->] (-0.3,0) -- (4.1,0) node[right] {$x$};
    \draw[->] (0,-0.3) -- (0,3.5) node[above] {$y$};
    \node[below left] at (0,0) {$O$};

    \node at (2,4.5) {$f^{-1}:[c,d]\longrightarrow[a,b]$};

    \draw[decorate,decoration={brace,amplitude=4pt,mirror}]
        (0.8,4.1) -- (3.2,4.1);
    \node at (2,3.7) {\small è una funzione};

    \draw plot[smooth] coordinates {
        (0.6,3) (1.1,2.9) (1.6,2.6)
        (2.1,2.0) (2.6,1.2) (3.2,0.5)
    };

    \fill (0.6,3) circle (2pt);
    \fill (3.2,0.5) circle (2pt);
    \fill (2.1,2.0) circle (2pt);

    \draw[dashed] (0.6,3) -- (0.6,0);
    \draw[dashed] (0,3) -- (0.6,3);
    \draw[dashed] (3.2,0.5) -- (3.2,0);
    \draw[dashed] (0,0.5) -- (3.2,0.5);

    \draw[->] (0,2.0) -- (2.1,2.0);
    \draw[->] (2.1,2.0) -- (2.1,0);

    \node[left] at (0,3) {$d$};
    \node[left] at (0,2.0) {$y_0$};
    \node[left] at (0,0.5) {$c$};

    \node[below] at (0.6,0) {$a$};
    \node[below] at (2.1,0) {$x_0$};
    \node[below] at (3.2,0) {$b$};

\end{tikzpicture}
\end{minipage}
\hfill
\begin{minipage}[t]{0.32\textwidth}
\vspace{0pt}
\centering
\begin{tikzpicture}[scale=0.78, >=Stealth, thick]

    \draw[->] (-0.3,0) -- (4.1,0) node[right] {$x$};
    \draw[->] (0,-0.3) -- (0,3.5) node[above] {$y$};
    \node[below left] at (0,0) {$O$};

    \node at (2,4.5) {$f^{-1}:[c,d]\longrightarrow[a,b]$};

    \draw[decorate,decoration={brace,amplitude=4pt,mirror}]
        (0.8,4.1) -- (3.2,4.1);
    \node at (2,3.7) {\small non è una funzione};

    \draw plot[smooth] coordinates {
        (0.6,3)
        (0.85,2.2)
        (1.1,1.5)
        (1.4,1.1)
        (1.7,1.15)
        (1.9,1.5)
        (2.2,1.85)
        (2.4,1.85)
        (2.6,1.5)
        (2.9,0.9)
        (3.2,0.5)
    };

    \fill (0.6,3) circle (2pt);
    \fill (3.2,0.5) circle (2pt);

    \draw[dashed] (0.6,3) -- (0.6,0);
    \draw[dashed] (0,3) -- (0.6,3);
    \draw[dashed] (3.2,0.5) -- (3.2,0);
    \draw[dashed] (0,0.5) -- (3.2,0.5);

    \draw[->] (0,1.5) -- (1.1,1.5);
    \draw[->] (1.1,1.5) -- (1.9,1.5);
    \draw[->] (1.9,1.5) -- (2.6,1.5);

    \draw[->] (1.1,1.5) -- (1.1,0);
    \draw[->] (1.9,1.5) -- (1.9,0);
    \draw[->] (2.6,1.5) -- (2.6,0);

    \node[left] at (0,3) {$d$};
    \node[left] at (0,1.5) {$y_0$};
    \node[left] at (0,0.5) {$c$};

    \node[below] at (0.6,0) {$a$};
    \node[below] at (1.1,0) {$x_0$};
    \node[below] at (1.9,0) {$x_1$};
    \node[below] at (2.6,0) {$x_2$};
    \node[below] at (3.2,0) {$b$};

\end{tikzpicture}
\end{minipage}
\par
\begin{definizione}[Funzione invertibile\footnote{Il simbolo $f^{-1}$ indica la funzione inversa di $f$, non ha il significato di $f$ elevato alla $-1$, cioè di $\dfrac{1}{f}$.}]
    Una funzione $f$ si dice invertibile se e solo se definisce una corrispondenza biunivoca tra il suo dominio e il suo insieme immagine: in tal caso, si chiama funzione inversa di $f$, e si indica con il simbolo $f^{-1}$, la funzione che associa a ciascun elemento di $f$ la sua controimmagine.
\end{definizione}
La condizione di invertibilità equivale a richiedere che la funzione sia iniettiva, cioè che ogni retta orizzontale intersechi il grafico in uno e un solo punto.

Si consideri $P(a,b)$ un punto appartenente al grafico della funzione $f$. Ciò significa che $f(a)=b$, vale a dire $f^{-1}(b)=a$. Dunque se $P(a,b)$ appartiene al grafico di $f$, allora $P'(b,a)$ appartiene al grafico di $f^{-1}$.

\begin{center}

\begin{tikzpicture}[scale=1.1, >=stealth]

    \draw[->] (-0.5,0) -- (4.5,0) node[right] {$x$};
    \draw[->] (0,-0.5) -- (0,4.5) node[above] {$y$};
    \node[below left] at (0,0) {$O$};

    \draw[thick, dashed] (-0.5,-0.5) -- (4.3,4.3);
    \node[above right] at (3.2,3.2) {$y=x$};

    \coordinate (O) at (0,0);
    \coordinate (P) at (3,1.5);
    \coordinate (Pp) at (1.5,3);

    \coordinate (A) at (3,0);
    \coordinate (B) at (0,3);

    \draw (O) -- (B) -- (Pp) -- cycle;
    \draw (O) -- (A) -- (P) -- cycle;

    \draw (Pp) -- (P);

    \draw[thick] (O) -- (P);
    \draw[thick] (O) -- (Pp);

    \draw (0,2.8) -- (0.2,2.8) -- (0.2,3);
    \draw (2.8,0) -- (2.8,0.2) -- (3,0.2);

    \draw (2.12,2.12) -- (2.26,1.98) -- (2.40,2.12);

    \node[above left] at (Pp) {$P'(b,a)$};
    \node[right] at (P) {$P(a,b)$};

\end{tikzpicture}

\end{center}
\begin{osservazione}
    Il grafico della funzione $f^{-1}$, inversa della funzione $f$, è il simmetrico del grafico di $f$ rispetto alla bisettrice del primo e del terzo quadrante.
\end{osservazione}
ESEMPIO:
\begin{center}
    \begin{tikzpicture}
  \begin{axis}[
    axis lines=middle,
    axis line style={-Stealth},
    axis equal,
    xmin=-6, xmax=8,
    ymin=-6, ymax=8,
    grid=none,
    major grid style={gray!25},
    minor tick num=0,
    xtick=\empty,
    ytick=\empty,
    xlabel={$x$},
    ylabel={$y$},
    legend style={
      at={(0.02,0.98)},
      anchor=north west,
      draw=none,
      fill=white,
      cells={align=left}
    },
    clip=false,
    every axis plot/.append style={line cap=round, line join=round}
  ]

    \addplot[
      black,
      dashed,
      thick,
      domain=-6:8,
      samples=2
    ] {x};
    \node[font=\small, anchor=west, inner sep=3pt] at (axis cs:8,8) {$y=x$};

    \addplot[
      black,
      very thick,
      mark=*,
      mark size=2.3pt
    ] coordinates {
      (-5,-5)
      (1,-2)
      (3,1)
      (7,3)
    };
    \node[font=\small, anchor=north, inner sep=2pt] at (axis cs:7,2.2) {$y=f(x)$};

    \foreach \x/\y/\pos in {
      -5/-5/below left,
       1/-2/below right,
       3/ 1/below right,
       7/ 3/above right
    }{
      \edef\tmp{\noexpand\node[font=\noexpand\footnotesize, inner sep=1pt] at (axis cs:\x,\y) [\pos] {$(\x,\y)$};}\tmp
    }

    \addplot[
      black,
      very thick,
      mark=*,
      mark size=2.3pt
    ] coordinates {
      (-5,-5)
      (-2,1)
      (1,3)
      (3,7)
    };
    \node[font=\small, anchor=east, inner sep=2pt] at (axis cs:2,6) {$y=f^{-1}(x)$};

    \foreach \x/\y/\pos in {
      -2/ 1/above left,
       1/ 3/above left,
       3/ 7/right
    }{
      \edef\tmp{\noexpand\node[font=\noexpand\footnotesize, inner sep=1pt] at (axis cs:\x,\y) [\pos] {$(\x,\y)$};}\tmp
    }

    

  \end{axis}
\end{tikzpicture}

\end{center}
Se $f$ è definita dall'equazione $y=f(x)$ allora $f^{-1}$ è definita dall'equazione: $x=f(y)$.

ESEMPIO: $f(x)=\sqrt[3]{x}-1$

$\begin{array}{l}
    y=\sqrt[3]{x}-1\\
    x=\sqrt[3]{y}-1\\
    \sqrt[3]{y}-1=x\\
    \sqrt[3]{y}=x+1\\
    y=(x+1)^3\\
    f^{-1}(x)=(x+1)^3
\end{array}$
\newpage
\begin{thebibliography}{9}
\bibitem{articolo1}
L. Sasso, E. Zoli,
\emph{Colori della Matematica, EDIZIONE VERDE, VOLUME 2},
DEASCUOLA, Petrini, 2022.

\bibitem{articolo2}
L. Sasso, E. Zoli,
\emph{Colori della Matematica, EDIZIONE VERDE, VOLUME 3},
DEASCUOLA, Petrini, 2022.
\end{thebibliography}
\end{document}